% Das Laborbuch: die logischen Schritte dieses Projekts, für Leserinnen und
% Leser — Kolleg:innen, Kooperationspartner:innen, man selbst in sechs Monaten. Vom Agenten
% in der Sprache der Forscherin / des Forschers geschrieben (labbook/voice.md),
% ein \section pro Schritt.
%
% Kästen, Abbildungen und Tabellen kommen aus dem Claim-Ledger über
% claims.gen.tex (`labbook.py gen`). Den Text eines Claims nie hierher
% kopieren: \claimbox{<id>} zeigt ihn genau so, wie das Ledger ihn hat.
\documentclass[11pt,a4paper]{article}
\usepackage[ngerman]{babel}
\usepackage{berry-labbook}
\renewcommand{\berryHoldsWhen}{Gilt, wenn}
\renewcommand{\berryConventions}{Konventionen}
\renewcommand{\berryRestsOn}{Beruht auf}
\renewcommand{\berryClaim}{Claim}
\renewcommand{\berryDraft}{Entwurf — noch nicht abgestimmt}
\renewcommand{\berrySealed}{versiegelt}
\renewcommand{\berryInLibrary}{in der Library}
\renewcommand{\berryFigureMissing}{Abbildung nicht erzeugt}
\renewcommand{\berryMissingClaim}{Kein Claim mit dieser Id im Ledger}
\renewcommand{\berryPaperReview}{Paper geprüft}
\renewcommand{\berryVerdictholds}{stimmt}
\renewcommand{\berryVerdictholdswithcorrection}{stimmt mit Korrektur}
\renewcommand{\berryVerdictfails}{falsch — nicht verwendet}
\renewcommand{\berryVerdictcouldnotcheck}{nicht prüfbar}
\renewcommand{\berryEvidence}{Evidenz}
\renewcommand{\berryRelcites}{zitiert}
\renewcommand{\berryReluses}{verwendet}
\renewcommand{\berryRelextends}{erweitert}
\renewcommand{\berryRelreplicates}{repliziert}
\renewcommand{\berryRelrefutes}{widerlegt}
\renewcommand{\berryTyperesult}{Ergebnis}
\renewcommand{\berryTypemethod}{Methode}
\renewcommand{\berryTypenegativeresult}{negatives Ergebnis}
\renewcommand{\berryTypecalibration}{Kalibrierung}
\input{claims.gen.tex}

\title{@TITLE@}
\author{@AUTHOR@}
\date{\today}

\begin{document}
\maketitle

\section*{Stand}
% Drei, vier Sätze: die Frage, was bisher gesichert ist, was offen ist.
% Wird neu geschrieben, wenn das Projekt weitergeht.

\berryoverview

% --- die Schritte --------------------------------------------------------------
%
% \section{<was dieser Schritt ergeben hat, in wenigen Worten>}
% Warum dieser Schritt: die Frage, und warum sie als nächste kam.
% Was gemacht wurde: die Methode in einfachen Worten; was festgehalten wurde.
% \claimbox{<Claim-Id>}
% \claimfigures{<Claim-Id>}
% Was das für den nächsten Schritt heißt. Geprüfte Paper: \sourcebox{<sources/…json>}

% --- Ende ----------------------------------------------------------------------

\section*{Offene Punkte und Folgefragen}
\begin{itemize}
  \item % Was nicht geklärt ist, was als Nächstes zu versuchen ist, was zu fragen ist.
\end{itemize}

% \section*{Notizen}
% Nur, was die Forscherin oder der Forscher über die Claims hinaus
% festgehalten haben will: Kontext, Entscheidungen, Anmerkungen. Nie anstelle
% eines Claims. Zum Verwenden das % entfernen.

\appendix
\section{Konventionen}
\berryconventions

\end{document}
